\documentclass{article}
% NeurIPS 2026 style (placeholder — download actual from neurips.cc)
\usepackage[preprint]{neurips_2026}
\usepackage{amsmath,amssymb,graphicx,booktabs,hyperref}

\title{Context Distraction in Small Language Models:\\
Why Retrieval-Augmented Generation Fails\\on Structured Code Generation Tasks}

\author{
William Chen \\
National Taiwan University \\
\texttt{william@ntu.edu.tw}
}

\begin{document}
\maketitle

\begin{abstract}
We present systematic evidence that retrieval-augmented generation (RAG)
and other context-augmentation strategies degrade the performance of small
language models ($<$10B parameters) on structured code generation tasks.
Through seven independent experiments on HLS pragma optimization—including
oracle experiments with perfect retrieval and classification—we demonstrate
that \emph{any} additional prompt context reduces pragma prediction accuracy
by 8–30\%. We term this phenomenon \textbf{Context Distraction}: small
models with limited attention capacity attend to surface patterns in
augmented context rather than reasoning about the target code structure.
Our findings are validated on 120 ForgeHLS kernels with Vitis HLS synthesis,
10 out-of-distribution custom kernels, and cross-architecture analysis
spanning DSP, cryptography, and image processing domains.
The only effective intervention is post-processing constraint integration
(+18.6\%), not pre-processing context augmentation.
These results establish fundamental boundaries for deploying small models
in domain-specific code generation and have implications for RAG system
design in resource-constrained environments.
\end{abstract}

% TODO: Expand each section to 1-2 pages for 9-page total
% Current content from paper2_context_distraction/main.tex
% Need to add:
% - Detailed experiment methodology
% - Per-kernel analysis
% - Scaling discussion
% - Attention analysis (if feasible)
% - Extended related work

\section{Introduction}
% Expand from current 0.5 page to 1.5 pages
% Add: motivation, problem statement, contributions list

\section{Background and Related Work}
% Expand from current 0.3 page to 1.5 pages
% Add: RAG theory, ICL in code generation, HLS optimization landscape

\section{Methodology}
% NEW section: describe PFS metric, post-process merge, experimental setup

\section{Experiments}
% Expand from current 1 page to 3 pages
% 4.1: Prompt strategy comparison (10 kernels)
% 4.2: Context length curve (5 kernels × 2 seeds)
% 4.3: Routing + Skills (30 kernels)
% 4.4: Oracle experiments (20 kernels)
% 4.5: Cross-architecture analysis (120 kernels)

\section{Analysis}
% 5.1: Why does context hurt? (attention competition)
% 5.2: Per-pragma-type breakdown
% 5.3: Implications for RAG system design

\section{Discussion and Limitations}
% Scale dependency, domain specificity, future work

\section{Conclusion}

\end{document}
